
\section{Preliminaries}
\label{sec:preliminary}

This section defines the problem of autonomous data preparation (ADP), describes the operators supported by \sys, and introduces background on agentic reasoning.

\subsection{Problem Formulation}
\label{subsec:problem_definition}
This paper studies the problem of \emph{autonomous data preparation} (ADP), which aims to transform and clean tabular data from multiple sources into a target table given only the specification of the target schema.
We formally define the problem as follows.

\stitle{Source and Target Data.}
We consider a set of \emph{source tables} $\mathcal{S} = \{S_1, \ldots, S_n\}$. Each source table is denoted as $S_i = (\Sigma_i, D_i)$, where $\Sigma_i$ denotes the schema of $S_i$ and $D_i$ denotes the set of tuples conforming to $\Sigma_i$. The schema $\Sigma_i$ captures both table-level and column-level specifications and is defined as $\Sigma_i = (\tau_i, C_i)$. Here, $\tau_i$ represents a textual description that summarizes the semantic meaning of the table, and $C_i = \{c_i^{1}, \ldots, c_i^{m_i}\}$ denotes the set of column specifications. Each column specification $c_i^{j}$ includes column-level metadata such as column name, data type, and semantic description.
In practice, the schema $\Sigma_i$ of a source table may be partially specified or entirely unavailable. In such cases, table- and column-level specifications can be inferred from the raw data in $D_i$. For example, each source table $S_i$ in Figure~\ref{fig:introduction_example} consists of a schema $\Sigma_i$ which may be specified by table- and column-level information, and a set $D_i$ of data tuples, \eg (1,``jurassic part'', 101, ``Sci-Fi'').
% For example, consider the ADP case in Figure~\ref{fig:introduction_example}. The schema $\Sigma_i$ of each source table $S_i$ in $\mathcal{S}$ is partially specified with table name (e.g., ``movies'') and column name (e.g., ``id''). The raw 

% \fanj{For example, \ldots}

Similarly, the \emph{target table} is denoted as $T=(\Sigma^{\ast}, D^{\ast})$, where $\Sigma^\ast = (\tau^\ast, C^\ast)$ specifies the target schema and $D^\ast$ denotes the target tuples to be produced. 
Note that, in autonomous data preparation, only the target schema $\Sigma^\ast$ is given as input, while the target tuples $D^\ast$ are unknown and must be generated by transforming and cleaning the source data.
For example, in Figure~\ref{fig:introduction_example}, the target schema $\Sigma^{\ast}$ is specified by a table-level description (\eg ``\emph{The table contains metrics \ldots}'') and a set of column specifications, each including a column name (\eg ``\emph{director\_name}'') and a corresponding semantic description (\eg ``\emph{The name of the director.}'').
% \fanj{For example, \ldots}

%Let $T$ be a source table consisting of a schema (columns) and instances (rows). Let $\mathbb{O}$ be the set of operators. We formulate the problem of Autonomous Data Preparation (ADP) as below.
%
%\begin{definition}
%    Given (1) a collection of \textit{source tables} $\mathcal{D}=\{T_1,\dots,T_M\}$ where $T$ denotes a source table consisting of a table name, a schema (columns) and instances (rows), and (2) \textit{a target specification} $\mathcal{S}=(\mathcal{S}_{table}, \mathcal{S}_{column})$ where $\mathcal{S}_{table}$ denotes a table-level summary of the target table and $\mathcal{S}_{column}$ denotes a structured constraint specifying the desired output schema. 
%    Synthesize a pipeline of parameterized operators $\Pi=[o_1, \dots, o_i]$, with $type(o_i) \in \mathbb{O}$, such that applying each operator $o_i$ successively on $\mathcal{D}$ with the operator execution engin $\mathcal{F}_{op}$ to produce a generated table $\hat{T}$ that aligns with the target specification $\mathcal{S}$.
%\end{definition}

\stitle{Data Prep Pipeline.}
We model the process of data prep as a composition of \emph{operators} applied to source tables.
Specifically, an operator represents a primitive data transformation. We consider a finite set of \emph{operator types}, where each type specifies the transformation semantics and the corresponding parameters. 

Formally, an operator $o$ is characterized by its type $\kappa$ and parameters $\theta$, and defines a function $o_{\theta}: \mathcal{T} \rightarrow \mathcal{T}$, where $\mathcal{T}$ denotes a set of tables, and each table is of the form $(\Sigma, D)$ as defined previously. Given an input table set $\mathcal{T}_{\tt in}$, the operator produces an output table set, \ie $\mathcal{T}_{\tt out}=o_{\theta}(\mathcal{T}_{\tt in})$.
Note that different operator types capture different classes of data transformations, such as filtering, joining, aggregation, and cleaning. The current operator types supported by \sys will be introduced in Section~\ref{subsec:operator_space}.

We then define a data prep pipeline (or pipeline for short) as a sequence of operators applied iteratively to transform source tables $\mathcal{S}$ into a target table $\hat{T}$. Formally, a pipeline $P$ is defined as an ordered composition of $k$ operators, \ie
%
\begin{equation} \nonumber
    P = o^{(k)}_{\theta_k} \circ \cdots \circ o^{(2)}_{\theta_2} \circ o^{(1)}_{\theta_1}.
\end{equation}
%
Given a set of source tables $\mathcal{S}$, the execution of pipeline $P$ produces a sequence of intermediate table sets, \ie 
%
\begin{equation} \nonumber
    \mathcal{T}_{0} = \mathcal{S}, ~\mathcal{T}_i =  o^{(i)}_{\theta_i}(\mathcal{T}_{i-1}), i=1,\ldots,k, 
\end{equation}
%
where $\mathcal{T}_k$ denotes the final output of the pipeline, which we refer to target table $\hat{T}$.
%
%\add{
Consider operator \texttt{Deduplicate (movies, [id], first)} as an example. 
Its operator type $\kappa$ is \texttt{Deduplicate} and its parameters $\theta$ refers to (\texttt{\textit{movies}, [id], first}). Executing this operator on source tables produces intermediate tables consisting of a cleaned \textit{movie}, while tables \textit{ratings} and \textit{directors} remain unchanged.
% For example, $\mathcal{T}_0$ is initialized with source tables illustrated in the left part of Figure~\ref{fig:introduction_example}. By executing operator, e.g., \texttt{Deduplicate} that removes duplicated record in table ``movies'' based on id values, we transform $\mathcal{T}_0$ to $\mathcal{T}_1$ which consists of updated table ``movies''
%}
% \fanj{For Example, \ldots}

\stitle{Problem Statement.}
Given a set of source tables $\mathcal{S} = \{S_1, \ldots, S_n\}$, a target schema $\Sigma^{\ast}$, and a set of available operator types $\mathcal{O}$, the autonomous data preparation (ADP) problem is to find an optimal data preparation pipeline $P$ that transforms $\mathcal{S}$ into a target table $T = (\Sigma^{\ast}, \hat{D})$, where $\hat{D}$ denotes the tuples produced by $P$ that satisfy the target schema $\Sigma^{\ast}$.
We assume a ground-truth target table $T^\ast=(\Sigma^\ast, D^\ast)$, which is used \emph{only} for evaluation. Thus, the optimality of a pipeline is defined by the quality of $\hat{D}$ with respect to $D^\ast$ under task-specific evaluation metrics, which will be introduced in the experimental setup (Section~\ref{sec:exp}).
%
\begin{example}

Figure~\ref{fig:introduction_example} provides an example of an ADP task with three source tables: \textit{movies}, \textit{ratings}, and \textit{directors}. The goal is to generate a target table that conforms to a given target schema. The target schema describes both the overall semantics of the table and its desired structure. Specifically, it provides a table-level summary of the intended content and specifies column-level requirements, including column names (e.g., \textit{genre}), semantic meanings (e.g., movie genre category), and data constraints (e.g., single-valued values).
To achieve this goal, we construct a data preparation pipeline composed of multiple operators, such as \texttt{Deduplicate} and \texttt{Join}. Each operator is applied to the current set of tables and produces updated intermediate tables. For example, applying \texttt{Join} merges the \textit{movies} and \textit{directors} tables on movie titles, yielding a table named \textit{movies\_directors\_join} that integrates movie data with director names, while keeping the \textit{ratings} table unchanged.
After executing all operators in the pipeline, the process produces the target table, shown on the right of Figure~\ref{fig:introduction_example}. 
%    Consider the ADP task $(\mathcal{D}, \mathcal{S})$ in Figure~\ref{fig:introduction_example}.
%    As shown in the left part, $\mathcal{D}$ contains three source tables which are named as \textit{movies}, \textit{ratings} and \textit{directors} respectively.
%    $\mathcal{S}=(\mathcal{S}_{table}, \mathcal{S}_{column})$ describes what does the target table look like. Specifically, $\mathcal{S}_{table}$ provides table-level summary, which is \textit{``The table contains''} and $\mathcal{S}_{column}$ provides description for the schema in the table which specifies the column name (e.g., ``genre''), the description for the semantic meanings (e.g., ``movie genre category'') and the requirement for the data (e.g., ``A single'').
%    It requires to generate a ground-truth operator pipeline $\Pi$ consisting of operators including \texttt{Deduplicate}, \texttt{Join}, etc.
%    These operators will be executed to modify source tables $\mathcal{D}$ by updating one or more tables in $\mathcal{D}$.
%    For example, we conduct \texttt{Deduplicate} to filter out duplicated rows in table \textit{movies} and update the source tables $\mathcal{D}$ with original table \textit{ratings}, original table \textit{directors} and a cleaned table \textit{movies}.
%    Finally, after executing all operators in the ground-truth operator pipeline, it outputs a table $\hat{T}$ shown on the right of Figure~\ref{fig:introduction_example}. The generated table contains data from source tables which aligns with the target specification $\mathcal{S}$.
\end{example}

\stitle{Remarks.}
The design of target schema $\Sigma^\ast$ removes the need for access to ground-truth target data, which is often unavailable in real-world settings. This distinguishes our problem from prior work~\cite{he2018transform, gulwani2012spreadsheet, jin2017foofah, singh2016blinkfill} that assumes access to target data examples.

%The design of the target specification $\mathcal{S}$ ensures that the method does not need the access to the data instances of ground-truth target table, which is difficulty to get in real application scenarios. This is the essential difference between our problem and the one that the previous work~\cite{he2018transform, gulwani2012spreadsheet, jin2017foofah, singh2016blinkfill} focused on.

\subsection{Supported Operator Types}
\label{subsec:operator_space}
\sys supports a wide range of data prep operators, which are summarized in Table~\ref{tbl:operator_space}. Note that the operator set is designed to be extensible, allowing new operator types to be incorporated as needed. 
chWe consider a total of 31 operators, which can be categorized into four groups. Due to space constraints, we present definitions of representative operators from each category, with the complete operator set described in~\cite{technical_report}.

\begin{table}[t]
  \centering
  \caption{\bf Overview of the Operator Space in \sys.}
  \vspace{-1em}
  \label{tbl:operator_space}
  \resizebox{1.0\columnwidth}{!}{
    \begin{tabular}{|c||l|}
      \hline
      \textbf{Category} & \textbf{Operators} \\
      \hline \hline

      % --- Cell-level ---
      \textbf{Cell-level}
      & \texttt{ValueTransform}, \texttt{StandardizeDatetime}, \texttt{CastType} \\
      \hline \hline

      % --- Row-level ---
      \multirow{3}{*}{\textbf{Row-level}}
      & \texttt{Filter}, \texttt{Sort}, \texttt{TopK}, \texttt{Deduplicate} \\
      & \texttt{DropNA}, \texttt{MissingValueImputation}, \texttt{ErrorDetection} \\
      & \texttt{OutlierDetection} \\
      \hline \hline

      % --- Column-level ---
      \multirow{2}{*}{\textbf{Column-level}}
      & \texttt{SplitColumn}, \texttt{RenameColumn}, \texttt{AddNewColumn}, \texttt{DropColumn} \\
      & \texttt{Concatenate}, \texttt{SelectColumn}, \texttt{Subtitle} \\
      \hline \hline

      % --- Table-level ---
      \multirow{2}{*}{\textbf{Table-level}}
      & \texttt{Pivot}, \texttt{Stack}, \texttt{WideToLong}, \texttt{Transpose}, \texttt{Explode} \\
      & \texttt{Join}, \texttt{Union}, \texttt{Append} \\
      \hline \hline

      % --- Aggregation ---
      \textbf{Aggregation}
      & \texttt{GroupBy}, \texttt{Count}, \texttt{CalculateStatistic} \\
      \hline \hline

      % --- Escape Hatch ---
      \textbf{Escape Hatch}
      & \texttt{ExeCode} \\
      \hline

      % If you must include Terminate, uncomment the following:
      % \hline \hline
      % \textbf{Control} & \texttt{Terminate} \\
      % \hline

    \end{tabular}
  }
  \vspace{-0.5em}
\end{table}




\add{
\begin{table*}[t]
  \centering
  \caption{\bf Taxonomy of supported data preparation operators by functional type.}
  \vspace{-0.6em}
  \label{tbl:operator_space_type}

  \small
  \setlength{\tabcolsep}{6pt}
  \renewcommand{\arraystretch}{1.12}

  \resizebox{\textwidth}{!}{
    \begin{tabular}{|l||l|}
      \hline
      \textbf{Operator Type} & \textbf{Operators} \\
      \hline \hline

      \textbf{Data Cleaning} &
      \texttt{MissingValueImputation},
      \texttt{Deduplicate},
      \texttt{ErrorDetection},
      \texttt{OutlierDetection},
      \texttt{DropNA} \\
      \hline

      \textbf{Value Normalization} &
      \texttt{ValueTransform},
      \texttt{StandardizeDatetime},
      \texttt{CastType} \\
      \hline

      \textbf{Schema Editing} &
      \texttt{RenameColumn},
      \texttt{SplitColumn},
      \texttt{AddNewColumn},
      \texttt{DropColumn},
      \texttt{Concatenate},
      \texttt{SelectColumn},
      \texttt{Subtitle} \\
      \hline

      \textbf{Row Selection} &
      \texttt{Filter},
      \texttt{Sort},
      \texttt{TopK} \\
      \hline

      \textbf{Aggregation} &
      \texttt{GroupBy},
      \texttt{Count},
      \texttt{CalculateStatistic} \\
      \hline

      \textbf{Table Combination} &
      \texttt{Join},
      \texttt{Union},
      \texttt{Append} \\
      \hline

      \textbf{Table Reshaping} &
      \texttt{Pivot},
      \texttt{Stack},
      \texttt{WideToLong},
      \texttt{Transpose},
      \texttt{Explode} \\
      \hline

      \textbf{Program Synthesis} &
      \texttt{ExeCode} \\
      \hline
    \end{tabular}
  }
\end{table*}

}



%Table~\ref{tbl:operator_space} illustrates an overview of 31 data preparation operators considered in this paper, which can be categorized into four groups. We provide definition of typical operators for each category due to space limit. A complete version can be found in~\cite{technical_report}.

\stitle{Data Cleaning Operators.} These operators address data quality issues such as inconsistencies, missing values, and outliers.
\begin{itemize}[leftmargin=1em]
    \item \texttt{Standardize(table, column, func)}. It unifies the format of a column based on user-defined transformation logic. For instance, to normalize the \texttt{title} column in Figure~\ref{fig:introduction_example}, the agent invokes this operator with \texttt{func} set to {\small\texttt{lambda x: ' '.join([s.capitalize() for s in x.split()])}}. This allows \sys to handle arbitrary formatting rules beyond simple casing.
    \item \texttt{DropNA(table, subset, how)}. It filters tuples containing null values. The \texttt{subset} argument restricts the inspection scope to specific columns, while \texttt{how} (``any'' or ``all'') determines the strictness of the elimination condition.
\end{itemize}

\stitle{Column Transformation.} These operators perform vertical manipulations on table attributes.
\begin{itemize}[leftmargin=1em]
    \item \texttt{SplitColumn(table, src\_col, func)}. Decomposes a composite column into multiple attributes using a Python function to parse complex patterns. For example, in Figure~\ref{fig:introduction_example}, splitting column `value' (e.g., ``8.5 (2021)'') into `Score' and `Year' is done via custom parsing logic in \texttt{func}.
    \item \texttt{RenameColumn(table, rename\_map)}. It takes a dictionary (e.g., {\small\texttt{\{'sid':'student\_id'\}}}) mapping current column names to their semantic counterparts in the target description.
\end{itemize}

\stitle{Table Reconstruction.} These operators modify the structural topology and row composition of the tables.
\begin{itemize}[leftmargin=1em]
    \item \texttt{Pivot(table, index, columns, values, aggfunc)}. Following~\cite{DBLP:journals/pvldb/LiHYWC23}, it restructures the table by rotating unique values from one column into multiple columns. The \texttt{aggfunc} parameter resolves cases where multiple source tuples map to a single cell.
    \item \texttt{Explode(table, column)}. It normalizes attributes containing list-like structures. For example, in Figure~\ref{fig:introduction_example}, this operator flattens the multi-valued \texttt{genres} column (``Action, Crime'') into distinct tuples, ensuring 1NF compliance.
\end{itemize}

\stitle{Auxiliary Operators.} These operator ensures the system to handle edge cases and complex logic beyond standard operators.
\begin{itemize}[leftmargin=1em]
    \item \texttt{CalculateStatistic(table, stat\_name, func)}. It derives scalar metadata (e.g., mean, variance) from a relation using aggregation logic encapsulated in \texttt{func}. The result is stored in a dedicated statistics registry for downstream decision-making.
    \item \texttt{ExeCode(tables, target, code)}. It serves as an ``escape hatch'' for requirements exceeding the expressivity of predefined operators. It allows the agent to synthesize arbitrary Python programs (e.g., invoking an external API or performing complex regex extraction) to generate a new table. This hybrid design integrates the reliability of atomic operators with the unbounded generalization of program synthesis.
\end{itemize}


\stitle{Comparison with Previous Work.} 
Our design distinguishes \sys from existing methods through three key advantages. 
First, in contrast to systems limited to isolated sub-tasks like table reconstruction~\cite{DBLP:journals/pvldb/LiHYWC23}, \sys provides a {unified scope} that simultaneously addresses cleaning, transformation, and reconstruction. 
Second, we mitigate the rigidity of static heuristic rules (e.g., Auto-Pipeline~\cite{yang2021auto}) by accepting executable Python logic as parameters. This enables the handling of complex, heterogeneous data patterns, e.g., regex-based splitting. 
Finally, the inclusion of a \texttt{ExeCode} operator ensures {extensibility}, serving as a reliable fallback for long-tail tasks that exceed the expressivity of standard operators.



\subsection{General Agentic Reasoning}
\label{subsec:agentic_reasoning_preliminary}

Agentic reasoning shifts from static text generation to autonomous problem solving. Here, the Large Language Model (LLM) acts as an autonomous {Agent}, interacting with an external execution environment to solve problems step by step. Formally, we model this reasoning as a tuple $\langle \mathcal{X}, \mathcal{A}, \mathcal{P}, \pi \rangle$:

\begin{itemize}[leftmargin=1em]
    \item {State Space ($\mathcal{X}$):} A state $x_t \in \mathcal{X}$ encapsulates the status of the task at time step $t$. For example, in code generation task, the status can be defined as all source files.
    
    \item {Action Space ($\mathcal{A}$):} At each time step, the agent selects a discrete action $a_t \in \mathcal{A}$. The action space encompasses both {cognitive actions} (e.g., planning, reasoning, analyzing) and {tool invocations} (e.g., invoking a Python interpreter to run Python code).
    
    \item {Environment Transition ($\mathcal{P}$):} The environment (i.e., the Python interpreter) takes the action $a_t$ and updates the system to a subsequent state $x_{t+1} = \mathcal{P}(x_t, a_t)$.
    
    \item {Agent Policy ($\pi$):} The agent acts according to a policy $\pi_\theta(a_t \mid x_t)$, whose parameters are defined by the LLM. The agent’s objective is to produce a \textit{trajectory} of actions $\tau = [a_0, a_1, \dots, a_L]$ that maximizes how well the final state aligns with the task objective.
\end{itemize}

\stitle{Connection with ADP.} This formulation is inherently tailored for Autonomous Data Preparation.
Specifically, ADP is inherently {state-dependent} where the validity of an operator (e.g., \texttt{SplitColumn}) depends on the actual data generated by the previous step. The transition function $\mathcal{P}$ enables the system to capture intermediate data snapshots and execution feedback (e.g., ``Column not found''), allowing the policy $\pi$ to perform dynamic error correction.
Thus, this formulation serves as the theoretical backbone of our proposed system \sys.



\iffalse

\begin{itemize}[leftmargin=*]
    \item \texttt{ValueTransform(table, column, func)}. It applies a string transformation function to a specific column. This is used to standardize string formats, such as removing special characters or fixing naming conventions.
    \item \texttt{StandardizeDatetime(table, column, format)}. It converts a column to a specific datetime format. It ensures temporal consistency across records by enforcing a target date format string.
    \item \texttt{CastType(table, column, dtype)}. It casts a column to a specified target data type (e.g., integer, float). This ensures the data schema aligns with the expected type requirements for downstream tasks.
    \item \texttt{MissingValueImputation(table, column, mode)}. It fills missing values in a column using a statistical method. The \texttt{mode} parameter specifies the strategy, such as using the mean, median, or mode of the column.
    \item \texttt{Filter(table, func)}. It selects rows based on a boolean condition defined by a custom function. The function takes a row (pd.Series) as input and returns \texttt{True} to keep the row or \texttt{False} to filter it out.
    \item \texttt{Sort(table, by, ascending)}. It orders rows based on specified columns and direction. The \texttt{ascending} parameter is a list of booleans corresponding to each column in the \texttt{by} list.
    \item \texttt{TopK(table, k)}. It retains only the top-$k$ rows of the table. This operator is typically used after a sorting operation to focus on the most significant records.
    \item \texttt{DropNA(table, subset, how)}. It removes rows containing null values. The \texttt{subset} parameter defines which columns to check, and \texttt{how} (``any'' or ``all'') controls the strictness of the removal.
    \item \texttt{Deduplicate(table, subset, keep)}. It removes duplicate rows based on a subset of columns. The \texttt{keep} parameter determines whether to retain the first or last occurrence of the duplicate entry.
    \item \texttt{ErrorDetection(table, column, func)}. It identifies invalid records using a custom verification function. Rows failing the validation logic (where the function returns \texttt{False}) are identified as errors.
    \item \texttt{OutlierDetection(table, column, action)}. It detects statistical outliers in a numerical column, typically using the IQR method. The \texttt{action} parameter dictates whether to delete these rows or tag them.
\end{itemize}

\subsubsection{Schema Editing ($\Sigma_i$)}
\begin{itemize}[leftmargin=1em]
    \item \texttt{RenameColumn(table, rename\_map)}. It takes a dictionary mapping current column names to their semantic counterparts. For example, \texttt{\{'sid': 'student\_id'\}} renames the ``sid'' column to ``student\_id''.
    \item \texttt{SplitColumn(table, source, target, func)}. It splits a source column into multiple target columns using a custom function. The function denotes the splitting logic in python function.
    \item \texttt{AddNewColumn(table, name, func)}. It creates a new column derived from existing data via a row-wise function. This is used to enrich the schema with computed attributes or derived features.
    \item \texttt{DropColumn(table, columns)}. It removes specified columns from the table. This reduces dimensionality by discarding irrelevant, redundant, or sensitive features.
    \item \texttt{Concatenate(table, columns, target, func)}. It merges multiple columns into a single string column using a formatting function. It is useful for combining fields like first and last names into a full name.
    \item \texttt{SelectColumn(table, columns)}. It retains only the specified list of columns, discarding all others. This focuses the table scope strictly on the relevant features needed for analysis.
    \item \texttt{Subtitle(table, title, target\_col)}. It adds a constant title or subtitle value as a new column to the table. This is often used for labeling datasets before merging or reporting.
\end{itemize}

\subsubsection{Table Reconstruction ($\Sigma_i + D_i$)}
\begin{itemize}[leftmargin=1em]
    \item \texttt{Pivot(table, index, columns, values, aggfunc)}. It reshapes data from long to wide format. It aggregates values using the specified \texttt{aggfunc} (e.g., ``sum'') based on the intersection of index and column identifiers.
    \item \texttt{Stack(table, id\_vars, value\_vars)}. It compresses multiple columns into a single key-value pair column (wide to long). This transforms column headers into row values, significantly increasing the table length.
    \item \texttt{WideToLong(table, stubnames, i, j)}. It converts wide-format panel data into long format by parsing column naming patterns. It uses prefixes (\texttt{stubnames}) and suffixes to extract variables and identifiers.
    \item \texttt{Transpose(table)}. It swaps the rows and columns of the table. This rotates the dataframe structure, turning indices into headers and vice versa.
    \item \texttt{Explode(table, column)}. It expands a column containing lists or arrays into multiple rows. Each element of the list becomes a separate row, replicating the data in other columns.
\end{itemize}

\subsubsection{Inter-table Composition \& Aggregation ($\mathcal{S}$)}
\begin{itemize}[leftmargin=1em]
    \item \texttt{Join(left, right, on, how)}. It merges two tables horizontally based on key columns. The \texttt{how} parameter specifies the join type (e.g., ``inner'', ``left'') to determine how unmatched rows are handled.
    \item \texttt{Union(tables, how)}. It combines multiple tables vertically into one. The \texttt{how} parameter determines if it is a ``union all'' (keeping duplicates) or ``union distinct''.
    \item \texttt{Append(table, other)}. It appends a specific table to the end of the source table. This is a specific vertical combination operation similar to union but strictly for appending one dataset to another.
    \item \texttt{GroupBy(table, by, agg)}. It groups data by specified columns and applies aggregation functions. The \texttt{agg} parameter defines a list of dictionaries specifying which column to aggregate and the function to use.
    \item \texttt{Count(table)}. It returns the total number of rows in the table. This provides a simple scalar metric representing the size of the dataset.
    \item \texttt{CalculateStatistic(table, stat, func)}. It computes a specific statistic (e.g., average, max) over the table using a custom function. It outputs a scalar value representing a global property of the dataset.
\end{itemize}

\subsubsection{Free-form Transformation ($\star$)}
\begin{itemize}[leftmargin=1em]
    \item \texttt{ExeCode(tables, target, func)}. It generates a new table from input tables using arbitrary Python code. This handles complex transformations that cannot be expressed by the standard set of operators.
\end{itemize}

\subsection{General Agentic Reasoning}
\label{subsec:agentic_reasoning_preliminary}

Agentic reasoning shifts from static text generation to autonomous problem solving. Here, the Large Language Model (LLM) acts as an autonomous {Agent}, interacting with an external execution environment to solve problems step by step. Formally, we model this reasoning as a tuple $\langle \mathcal{X}, \mathcal{A}, \mathcal{P}, \pi \rangle$:

\begin{itemize}[leftmargin=1em]
    \item {State Space ($\mathcal{X}$):} A state $x_t \in \mathcal{X}$ encapsulates the status of the task at time step $t$. For example, in code generation task, the status can be defined as all source files.
    
    \item {Action Space ($\mathcal{A}$):} At each time step, the agent selects a discrete action $a_t \in \mathcal{A}$. The action space encompasses both {cognitive actions} (e.g., planning, reasoning, analyzing) and {tool invocations} (e.g., invoking a Python interpreter to run Python code).
    
    \item {Environment Transition ($\mathcal{P}$):} The environment (i.e., the Python interpreter) takes the action $a_t$ and updates the system to a subsequent state $x_{t+1} = \mathcal{P}(x_t, a_t)$.
    
    \item {Agent Policy ($\pi$):} The agent acts according to a policy $\pi_\theta(a_t \mid x_t)$, whose parameters are defined by the LLM. The agent’s objective is to produce a \textit{trajectory} of actions $\tau = [a_0, a_1, \dots, a_L]$ that maximizes how well the final state aligns with the task objective.
\end{itemize}

\stitle{Connection with ADP.} This formulation is inherently tailored for Autonomous Data Preparation.
Specifically, ADP is inherently {state-dependent} where the validity of an operator (e.g., \texttt{SplitColumn}) depends on the actual data generated by the previous step. The transition function $\mathcal{P}$ enables the system to capture intermediate data snapshots and execution feedback (e.g., ``Column not found''), allowing the policy $\pi$ to perform dynamic error correction.
Thus, this formulation serves as the theoretical backbone of our proposed system \sys.

\fi